% ====================================================================
% ФАЙЛ: tasks/01_mechanics/friction_force_bank.tex
% ТЕМА: СИЛА ТРЕНИЯ
% ПАЧКА 1: ВАРИАНТЫ 1–10
% ====================================================================

% === ЗАДАЧА 1 ===
% Тег: #МЕХАНИКА #КИМ_05 #ВАРИАНТ_03
\begin{taskbasic}[code={\label{task:friction_v3_n5}}]
\textbf{Задача 1.} В лабораторной работе изучали движение небольшого бруска массой $200$~г по горизонтальной шероховатой поверхности под действием горизонтальной постоянной силы, равной по модулю $800$~мН. Зависимость модуля скорости бруска от времени приведена в таблице.
\begin{center}
\begin{tabular}{|c|c|c|c|c|c|c|c|}
\hline
Время $t$, с & 0 & 0,1 & 0,2 & 0,3 & 0,4 & 0,5 & 0,6 \\ \hline
Скорость $v$, м/с & 0 & 0,2 & 0,4 & 0,6 & 0,8 & 1,0 & 1,2 \\ \hline
\end{tabular}
\end{center}
Выберите все верные утверждения на основании анализа представленной таблицы.

1) Модуль ускорения бруска равен $2$~м/с$^{2}$.
2) Коэффициент трения бруска о поверхность $\mu = 0{,}3$.
3) Брусок движется равноускоренно.
4) В момент времени $0{,}3$~с кинетическая энергия бруска равна $36$~мДж.
5) Модуль равнодействующей сил, действующих на брусок, равен $0{,}8$~Н.

\kim{5}
\end{taskbasic}
\answer{134}

% === ЗАДАЧА 2 ===
% Тег: #МЕХАНИКА #КИМ_05 #ВАРИАНТ_04
\begin{taskbasic}[code={\label{task:friction_v4_n5}}]
\textbf{Задача 2.} В лабораторной работе изучали движение небольшого бруска массой $200$~г по горизонтальной шероховатой поверхности под действием горизонтальной постоянной силы, равной по модулю $800$~мН. Зависимость модуля скорости бруска от времени приведена в таблице.
\begin{center}
\begin{tabular}{|c|c|c|c|c|c|c|c|}
\hline
Время $t$, с & 0 & 0,1 & 0,2 & 0,3 & 0,4 & 0,5 & 0,6 \\ \hline
Скорость $v$, м/с & 0 & 0,2 & 0,4 & 0,6 & 0,8 & 1,0 & 1,2 \\ \hline
\end{tabular}
\end{center}
Выберите все верные утверждения на основании анализа представленной таблицы.

1) Модуль ускорения бруска равен $0{,}4$~м/с$^{2}$.
2) Коэффициент трения бруска о поверхность $\mu = 0{,}2$.
3) Брусок движется равномерно.
4) В момент времени $0{,}3$~с кинетическая энергия бруска равна $18$~мДж.
5) Модуль равнодействующей сил, действующих на брусок, равен $0{,}4$~Н.

\kim{5}
\end{taskbasic}
\answer{25}

% ====================================================================
% ФАЙЛ: tasks/01_mechanics/friction_force_bank_part2.tex
% ТЕМА: СИЛА ТРЕНИЯ (ДИНАМИКА СВЯЗАННЫХ ТЕЛ)
% ПАЧКА 2: ВАРИАНТЫ 11–20
% ====================================================================

% === ЗАДАЧА 1 ===
% Тег: #МЕХАНИКА #КИМ_26 #ВАРИАНТ_11
\begin{taskbasic}[code={\label{task:friction_v11_n26}}]
\textbf{Задача 1.} Если в установке первоначально покоящуюся тележку толкнуть влево, то она движется с ускорением $3$~м/с$^{2}$. Если же тележку толкнуть вправо, то она движется равномерно. Найдите массу $M$ тележки, если масса грузика на нити $m=150$~г. Массами блока и нити пренебречь. Нить нерастяжима. Модуль силы сопротивления движению тележки считать постоянным и одинаковым в обоих случаях, трением в оси блока пренебречь.

\nopagebreak\smallskip
\begin{center}
\begin{tikzpicture}[scale=1.2, every node/.style={font=\footnotesize}]
    % Стол
    \draw[thick, fill=gray!5] (-2.5,0) rectangle (3.0,0.15);
    % Тележка M
    \draw[thick, fill=brandPrimary!10] (-0.8,0.22) rectangle (-0.2,0.42);
    \node[above] at (-0.5,0.42) {$M$};
    \draw[thick, fill=black!50] (-0.65,0.185) circle (0.04);
    \draw[thick, fill=black!50] (-0.35,0.185) circle (0.04);
    % Блок
    \draw[thick, fill=gray!40] (3.0,0.15) circle (0.12);
    \fill[black] (3.0,0.15) circle (1pt);
    \draw[thick] (2.9,0.05) -- (3.0,0.15);
    % Груз m
    \draw[thick, fill=brandAccent!20] (2.85,-0.9) rectangle (3.15,-0.83);
    \node[right] at (3.2,-0.86) {$m$};
    % Нити
    \draw[thick] (-0.2,0.32) -- (3.0,0.27);
    \draw[thick] (3.12,0.15) -- (3.12,-0.83);
\end{tikzpicture}
\end{center}

\kim{26}
\end{taskbasic}
\answer{0{,}85~кг}

% === ЗАДАЧА 2 ===
% Тег: #МЕХАНИКА #КИМ_26 #ВАРИАНТ_12
\begin{taskbasic}[code={\label{task:friction_v12_n26}}]
\textbf{Задача 2.} Если в установке первоначально покоящуюся тележку толкнуть влево, то она движется с ускорением $2$~м/с$^{2}$. Если же тележку толкнуть вправо, то она движется равномерно. Найдите массу $m$ грузика на нити, если масса тележки $M=450$~г. Массами блока и нити пренебречь. Нить нерастяжима. Модуль силы сопротивления движению тележки считать постоянным и одинаковым в обоих случаях, трением в оси блока пренебречь.

\nopagebreak\smallskip
\begin{center}
\begin{tikzpicture}[scale=1.2, every node/.style={font=\footnotesize}]
    % Стол
    \draw[thick, fill=gray!5] (-2.5,0) rectangle (3.0,0.15);
    % Тележка M
    \draw[thick, fill=brandPrimary!10] (-0.8,0.22) rectangle (-0.2,0.42);
    \node[above] at (-0.5,0.42) {$M$};
    \draw[thick, fill=black!50] (-0.65,0.185) circle (0.04);
    \draw[thick, fill=black!50] (-0.35,0.185) circle (0.04);
    % Блок
    \draw[thick, fill=gray!40] (3.0,0.15) circle (0.12);
    \fill[black] (3.0,0.15) circle (1pt);
    \draw[thick] (2.9,0.05) -- (3.0,0.15);
    % Груз m
    \draw[thick, fill=brandAccent!20] (2.85,-0.9) rectangle (3.15,-0.83);
    \node[right] at (3.2,-0.86) {$m$};
    % Нити
    \draw[thick] (-0.2,0.32) -- (3.0,0.27);
    \draw[thick] (3.12,0.15) -- (3.12,-0.83);
\end{tikzpicture}
\end{center}

\kim{26}
\end{taskbasic}
\answer{0{,}045~кг}

% === ЗАДАЧА 3 ===
% Тег: #МЕХАНИКА #КИМ_26 #ВАРИАНТ_13
\begin{taskbasic}[code={\label{task:friction_v13_n26}}]
\textbf{Задача 3.} На горизонтальном столе находится брусок массой $M=1$~кг, соединённый невесомой нерастяжимой нитью, перекинутой через гладкий невесомый блок, с грузом массой $m=500$~г. На брусок действует сила величиной $F=9$~Н, направленная под углом $\alpha=30^{\circ}$ к горизонту. В момент начала движения груз находится на расстоянии $L=40$~см от края стола. За какое время груз поднимется до края стола, если коэффициент трения между бруском и столом $\mu=0{,}3$? Трением в оси блока и трением о воздух пренебречь.

\nopagebreak\smallskip
\begin{center}
\begin{tikzpicture}[scale=1.2, every node/.style={font=\footnotesize}]
    % Поверхность стола и угловой обрыв (чистый fill=gray!5)
    \draw[thick, fill=gray!5] (-1.5,0) -- (3.5,0) -- (3.5,-0.1) -- (-1.4,-0.1) -- (-1.4,-2.2) -- (-1.5,-2.2) -- cycle;
    % Штриховка обрыва стола
    \foreach \y in {-0.2,-0.4,-0.6,-0.8,-1.0,-1.2,-1.4,-1.6,-1.8,-2.0}
        \draw[thin] (-1.5,\y) -- (-1.4,\y+0.1);

    % Неподвижный блок на углу стола
    \draw[thick, fill=gray!30] (-1.5,0) circle (0.18);
    \draw[thick, fill=gray!10] (-1.5,0) circle (0.06);
    \fill[black] (-1.5,0) circle (1.2pt);

    % Брусок M (на столе)
    \draw[thick, fill=brandPrimary!10] (0.5,0) rectangle (1.7,0.7);
    \node at (1.1,0.35) {$M$};
    
    % Сила F, приложенная к бруску M
    \draw[-{Stealth[scale=0.8]}, ultra thick, brandAccent] (1.7,0.35) -- (2.7,0.85) node[above] {$\vec{F}$};
    \draw[dashed, thick] (1.7,0.35) -- (2.7,0.35);
    \draw (2.1,0.35) arc (0:26.5:0.4);
    \node at (2.25,0.47) {$\alpha$};
    
    % Груз m (висит вертикально)
    \draw[thick, fill=brandAccent!20] (-1.9,-1.5) rectangle (-1.45,-1.0);
    \node at (-1.675,-1.25) {$m$};
    
    % Горизонтальный участок нити (от блока к бруску M)
    \draw[thick] (-1.5,0.18) -- (0.5,0.18);
    
    % Вертикальный участок нити (от блока к грузу m)
    \draw[thick] (-1.68,0) -- (-1.68,-1.0);
    
    % Выносные линии и стрелка для расстояния L
    \draw[gray] (-1.4,0) -- (-1.0,0);
    \draw[gray] (-1.45,-1.0) -- (-1.0,-1.0);
    \draw[<->, gray] (-1.1,0) -- (-1.1,-1.0) node[midway, right, black] {$L$};
\end{tikzpicture}
\end{center}

\kim{26}
\end{taskbasic}
\answer{1~с}

% === ЗАДАЧА 4 ===
% Тег: #МЕХАНИКА #КИМ_26 #ВАРИАНТ_14
\begin{taskbasic}[code={\label{task:friction_v14_n26}}]
\textbf{Задача 4.} На горизонтальном столе находится брусок массой $M=1$~кг, соединённый невесомой нерастяжимой нитью, перекинутой через гладкий невесомый блок, с грузом массой $m=500$~г. На брусок действует сила величиной $F=8$~Н, направленная под углом $\alpha=30^{\circ}$ к горизонту. В момент начала движения груз находится на расстоянии $L=1$~м от края стола. Каков коэффициент трения между бруском и столом, если груз поднимается до края стола за время $t=2$~с? Трением в оси блока и трением о воздух пренебречь.

\nopagebreak\smallskip
\begin{center}
\begin{tikzpicture}[scale=1.2, every node/.style={font=\footnotesize}]
    % Поверхность стола и угловой обрыв (чистый fill=gray!5)
    \draw[thick, fill=gray!5] (-1.5,0) -- (3.5,0) -- (3.5,-0.1) -- (-1.4,-0.1) -- (-1.4,-2.2) -- (-1.5,-2.2) -- cycle;
    % Штриховка обрыва стола
    \foreach \y in {-0.2,-0.4,-0.6,-0.8,-1.0,-1.2,-1.4,-1.6,-1.8,-2.0}
        \draw[thin] (-1.5,\y) -- (-1.4,\y+0.1);

    % Неподвижный блок на углу стола
    \draw[thick, fill=gray!30] (-1.5,0) circle (0.18);
    \draw[thick, fill=gray!10] (-1.5,0) circle (0.06);
    \fill[black] (-1.5,0) circle (1.2pt);

    % Брусок M (на столе)
    \draw[thick, fill=brandPrimary!10] (0.5,0) rectangle (1.7,0.7);
    \node at (1.1,0.35) {$M$};
    
    % Сила F, приложенная к бруску M
    \draw[-{Stealth[scale=0.8]}, ultra thick, brandAccent] (1.7,0.35) -- (2.7,0.85) node[above] {$\vec{F}$};
    \draw[dashed, thick] (1.7,0.35) -- (2.7,0.35);
    \draw (2.1,0.35) arc (0:26.5:0.4);
    \node at (2.25,0.47) {$\alpha$};
    
    % Груз m (висит вертикально)
    \draw[thick, fill=brandAccent!20] (-1.9,-1.5) rectangle (-1.45,-1.0);
    \node at (-1.675,-1.25) {$m$};
    
    % Горизонтальный участок нити (от блока к бруску M)
    \draw[thick] (-1.5,0.18) -- (0.5,0.18);
    
    % Вертикальный участок нити (от блока к грузу m)
    \draw[thick] (-1.68,0) -- (-1.68,-1.0);
    
    % Выносные линии и стрелка для расстояния L
    \draw[gray] (-1.4,0) -- (-1.0,0);
    \draw[gray] (-1.45,-1.0) -- (-1.0,-1.0);
    \draw[<->, gray] (-1.1,0) -- (-1.1,-1.0) node[midway, right, black] {$L$};
\end{tikzpicture}
\end{center}

\kim{26}
\end{taskbasic}
\answer{0{,}2}

% === ЗАДАЧА 5 ===
% Тег: #МЕХАНИКА #КИМ_26 #ВАРИАНТ_17
\begin{taskbasic}[code={\label{task:friction_v17_n26}}]
\textbf{Задача 5.} На гладком горизонтальном столе лежит доска массой $M=1$~кг и длиной $L=50$~см. На левом краю доски находится маленький брусок массой $m=200$~г. Брусок и доска связаны невесомой нерастяжимой нитью, перекинутой через невесомый гладкий блок, закреплённый на стене (отрезки нити, не лежащие на блоке, горизонтальны и параллельны). Коэффициент трения между бруском и доской $\mu=0{,}3$. Брусок начинают тянуть вправо постоянной силой $\vec{F}$, параллельной горизонтальным отрезкам нити. Через время $t=2$~с после начала движения брусок соскальзывает с доски. Определите модуль силы $\vec{F}$. Размерами бруска пренебречь.

\nopagebreak\smallskip
\begin{center}
\begin{tikzpicture}[scale=1.2, every node/.style={font=\footnotesize}]
    % Поверхность пола (горизонтальная линия стола)
    \draw[thick, fill=gray!5] (-2.2,0) -- (4.2,0) -- (4.2,-0.1) -- (-2.2,-0.1) -- cycle;
    
    % Левая вертикальная стена (опора для блока)
    \draw[thick, fill=gray!20] (-2.2,0) rectangle (-2.0,2.0);
    % Штриховка стены
    \foreach \y in {0.1,0.3,0.5,0.7,0.9,1.1,1.3,1.5,1.7,1.9}
        \draw[thin] (-2.2,\y) -- (-2.0,\y-0.1);

    % Доска M (большая нижняя доска)
    \draw[thick, fill=brandPrimary!10] (1.0,0) rectangle (3.4,0.5);
    \node at (2.2,0.25) {$M$};
    
    % Брусок m (маленький верхний брусок, на левом краю доски M)
    \draw[thick, fill=brandAccent!20] (1.0,0.5) rectangle (1.7,1.0);
    \node at (1.35,0.75) {$m$};
    
    % Неподвижный блок на левой стене
    \draw[thick, fill=gray!30] (-1.4,0.65) circle (0.25);
    \draw[thick, fill=gray!10] (-1.4,0.65) circle (0.08);
    \fill[black] (-1.4,0.65) circle (1.5pt);
    
    % Крепление блока к левой стене
    \draw[thick] (-2.0,0.65) -- (-1.4,0.65);
    
    % Верхняя нить (от блока к бруску m)
    \draw[thick] (-1.4,0.9) -- (1.0,0.9);
    
    % Нижняя нить (от блока к доске M)
    \draw[thick] (-1.4,0.4) -- (1.0,0.4);
    
    % Сила F, тянущая брусок m вправо
    \draw[-{Stealth[scale=0.8]}, ultra thick, brandAccent] (1.7,0.75) -- (3.0,0.75) node[above] {$\vec{F}$};
    
    % Выносные линии и стрелка для длины L доски M
    \draw[gray] (1.0,0) -- (1.0,-0.4);
    \draw[gray] (3.4,0) -- (3.4,-0.4);
    \draw[<->, gray] (1.0,-0.3) -- (3.4,-0.3) node[midway, below, black] {$L$};
\end{tikzpicture}
\end{center}

\kim{26}
\end{taskbasic}
\answer{1{,}35~Н}

% === ЗАДАЧА 6 ===
% Тег: #МЕХАНИКА #КИМ_26 #ВАРИАНТ_18
\begin{taskbasic}[code={\label{task:friction_v18_n26}}]
\textbf{Задача 6.} На гладком горизонтальном столе лежит доска массой $M=1$~кг и длиной $L=50$~см. На левом краю доски находится маленький брусок массой $m=200$~г. Брусок и доска связаны невесомой нерастяжимой нитью, перекинутой через невесомый гладкий блок, закреплённый на стене (отрезки нити, не лежащие на блоке, горизонтальны и параллельны). Коэффициент трения между бруском и доской $\mu=0{,}2$. Брусок начинают тянуть вправо с постоянной силой $\vec{F}$, параллельной горизонтальным отрезкам нити. Модуль этой силы $F=1{,}2$~Н. Через какое время после начала движения брусок соскользнёт с доски? Размерами бруска пренебречь.

\nopagebreak\smallskip
\begin{center}
\begin{tikzpicture}[scale=1.2, every node/.style={font=\footnotesize}]
    % Поверхность пола (горизонтальная линия стола)
    \draw[thick, fill=gray!5] (-2.2,0) -- (4.2,0) -- (4.2,-0.1) -- (-2.2,-0.1) -- cycle;
    
    % Левая вертикальная стена (опора для блока)
    \draw[thick, fill=gray!20] (-2.2,0) rectangle (-2.0,2.0);
    % Штриховка стены
    \foreach \y in {0.1,0.3,0.5,0.7,0.9,1.1,1.3,1.5,1.7,1.9}
        \draw[thin] (-2.2,\y) -- (-2.0,\y-0.1);

    % Доска M (большая нижняя доска)
    \draw[thick, fill=brandPrimary!10] (1.0,0) rectangle (3.4,0.5);
    \node at (2.2,0.25) {$M$};
    
    % Брусок m (маленький верхний брусок, на левом краю доски M)
    \draw[thick, fill=brandAccent!20] (1.0,0.5) rectangle (1.7,1.0);
    \node at (1.35,0.75) {$m$};
    
    % Неподвижный блок на левой стене
    \draw[thick, fill=gray!30] (-1.4,0.65) circle (0.25);
    \draw[thick, fill=gray!10] (-1.4,0.65) circle (0.08);
    \fill[black] (-1.4,0.65) circle (1.5pt);
    
    % Крепление блока к левой стене
    \draw[thick] (-2.0,0.65) -- (-1.4,0.65);
    
    % Верхняя нить (от блока к бруску m)
    \draw[thick] (-1.4,0.9) -- (1.0,0.9);
    
    % Нижняя нить (от блока к доске M)
    \draw[thick] (-1.4,0.4) -- (1.0,0.4);
    
    % Сила F, тянущая брусок m вправо
    \draw[-{Stealth[scale=0.8]}, ultra thick, brandAccent] (1.7,0.75) -- (3.0,0.75) node[above] {$\vec{F}$};
    
    % Выносные линии и стрелка для длины L доски M
    \draw[gray] (1.0,0) -- (1.0,-0.4);
    \draw[gray] (3.4,0) -- (3.4,-0.4);
    \draw[<->, gray] (1.0,-0.3) -- (3.4,-0.3) node[midway, below, black] {$L$};
\end{tikzpicture}
\end{center}

\kim{26}
\end{taskbasic}
\answer{1{,}22~с}

% === ЗАДАЧА 7 ===
% Тег: #МЕХАНИКА #КИМ_26 #ВАРИАНТ_19
\begin{taskbasic}[code={\label{task:friction_v19_n26}}]
\textbf{Задача 7.} На горизонтальном неподвижном столе лежит доска, на которой находится маленький брусок массой $m=400$~г. Брусок и доска связаны невесомой нерастяжимой нитью, перекинутой через невесомый блок, закреплённый на стене (отрезки нити, не лежащие на блоке, горизонтальны и параллельны). Коэффициент трения между бруском и доской $\mu_1=0{,}5$, между столом и доской $\mu_2=0{,}3$. Доску тянут вправо постоянной горизонтальной силой $F$, параллельной горизонтальным отрезкам нити, $F=12$~Н. Чему равна масса доски $M$, если модуль ускорения бруска относительно стола $a=2$~м/с$^{2}$? Трением в оси блока пренебречь.

\nopagebreak\smallskip
\begin{center}
\begin{tikzpicture}[scale=1.2, every node/.style={font=\footnotesize}]
    % Поверхность пола (горизонтальная линия стола)
    \draw[thick, fill=gray!5] (-2.2,0) -- (4.2,0) -- (4.2,-0.1) -- (-2.2,-0.1) -- cycle;
    
    % Левая вертикальная стена (опора для блока)
    \draw[thick, fill=gray!20] (-2.2,0) rectangle (-2.0,2.0);
    % Штриховка стены
    \foreach \y in {0.1,0.3,0.5,0.7,0.9,1.1,1.3,1.5,1.7,1.9}
        \draw[thin] (-2.2,\y) -- (-2.0,\y-0.1);

    % Доска M (большая нижняя доска)
    \draw[thick, fill=brandPrimary!10] (0.2,0) rectangle (2.6,0.5);
    \node at (1.4,0.25) {$M$};
    
    % Брусок m (маленький верхний брусок)
    \draw[thick, fill=brandAccent!20] (1.9,0.5) rectangle (2.6,1.0);
    \node at (2.25,0.75) {$m$};
    
    % Неподвижный блок на левой стене
    \draw[thick, fill=gray!30] (-1.4,0.65) circle (0.25);
    \draw[thick, fill=gray!10] (-1.4,0.65) circle (0.08);
    \fill[black] (-1.4,0.65) circle (1.5pt);
    
    % Крепление блока к левой стене
    \draw[thick] (-2.0,0.65) -- (-1.4,0.65);
    
    % Верхняя нить (от блока к бруску m)
    \draw[thick] (-1.4,0.9) -- (1.9,0.9);
    
    % Нижняя нить (от блока к доске M)
    \draw[thick] (-1.4,0.4) -- (0.2,0.4);
    
    % Сила F, тянущая доску M вправо
    \draw[-{Stealth[scale=0.8]}, ultra thick, brandPrimary] (2.6,0.25) -- (3.6,0.25) node[above] {$\vec{F}$};
\end{tikzpicture}
\end{center}

\kim{26}
\end{taskbasic}
\answer{1{,}2~кг}

% === ЗАДАЧА 8 ===
% Тег: #МЕХАНИКА #КИМ_26 #ВАРИАНТ_20
\begin{taskbasic}[code={\label{task:friction_v20_n26}}]
\textbf{Задача 8.} На горизонтальном неподвижном столе лежит доска массой $M=0{,}8$~кг. На доске находится маленький брусок массой $m=200$~г. Брусок и доска связаны невесомой нерастяжимой нитью, перекинутой через невесомый блок, закреплённый на стене (отрезки нити, не лежащие на блоке, горизонтальны и параллельны). Коэффициент трения между бруском и доской $\mu_1=0{,}5$, между столом и доской $\mu_2=0{,}3$. Доску тянут вправо постоянной силой $\vec{F}$, параллельной горизонтальным отрезкам нити. Чему равен модуль силы $\vec{F}$, если модуль ускорения бруска относительно стола $a=1$~м/с$^{2}$? Трением в оси блока пренебречь.

\nopagebreak\smallskip
\begin{center}
\begin{tikzpicture}[scale=1.2, every node/.style={font=\footnotesize}]
    % Поверхность пола (горизонтальная линия стола)
    \draw[thick, fill=gray!5] (-2.2,0) -- (4.2,0) -- (4.2,-0.1) -- (-2.2,-0.1) -- cycle;
    
    % Левая вертикальная стена (опора для блока)
    \draw[thick, fill=gray!20] (-2.2,0) rectangle (-2.0,2.0);
    % Штриховка стены
    \foreach \y in {0.1,0.3,0.5,0.7,0.9,1.1,1.3,1.5,1.7,1.9}
        \draw[thin] (-2.2,\y) -- (-2.0,\y-0.1);

    % Доска M (большая нижняя доска)
    \draw[thick, fill=brandPrimary!10] (0.2,0) rectangle (2.6,0.5);
    \node at (1.4,0.25) {$M$};
    
    % Брусок m (маленький верхний брусок)
    \draw[thick, fill=brandAccent!20] (1.9,0.5) rectangle (2.6,1.0);
    \node at (2.25,0.75) {$m$};
    
    % Неподвижный блок на левой стене
    \draw[thick, fill=gray!30] (-1.4,0.65) circle (0.25);
    \draw[thick, fill=gray!10] (-1.4,0.65) circle (0.08);
    \fill[black] (-1.4,0.65) circle (1.5pt);
    
    % Крепление блока к левой стене
    \draw[thick] (-2.0,0.65) -- (-1.4,0.65);
    
    % Верхняя нить (от блока к бруску m)
    \draw[thick] (-1.4,0.9) -- (1.9,0.9);
    
    % Нижняя нить (от блока к доске M)
    \draw[thick] (-1.4,0.4) -- (0.2,0.4);
    
    % Сила F, тянущая доску M вправо
    \draw[-{Stealth[scale=0.8]}, ultra thick, brandPrimary] (2.6,0.25) -- (3.6,0.25) node[above] {$\vec{F}$};
\end{tikzpicture}
\end{center}

\kim{26}
\end{taskbasic}
\answer{6~Н}

